\documentclass[11pt,a4paper]{article}
\usepackage[T1]{fontenc}
\usepackage[utf8]{inputenc}
\usepackage{lmodern}
\usepackage[a4paper,margin=2.7cm]{geometry}
\usepackage{microtype}
\usepackage{amsmath,amssymb}
\usepackage{setspace}
\usepackage{xcolor}
\usepackage{hyperref}
\hypersetup{colorlinks=true,linkcolor=black,urlcolor=blue}
\setlength{\parskip}{0.7em}
\setlength{\parindent}{0pt}
\onehalfspacing

\title{\textbf{The Elephant and the Agreement}\\
\large A Parable of Autonomy, Interdependence, and the Search for Truth}
\author{Albert Jan van Hoek}
\date{September 2026}

\begin{document}
\maketitle

\begin{quote}
\emph{None of them could see the whole elephant. That did not make their perceptions worthless. It made their relationship important.}
\end{quote}

\section*{The elephant arrives}

One morning, an elephant arrived in a village where six blind people lived.

None of them had ever encountered an elephant before. They had heard the word, but the word did not tell them what an elephant was.

So they went to meet it.

The first person reached the elephant's side. He placed both hands against the animal's broad flank.

``An elephant,'' he said, ``is like a wall.''

The second found a leg. She wrapped her arms around it.

``No,'' she said. ``An elephant is like a tree.''

The third found the trunk. It curled and moved beneath his hands.

``You are both mistaken. An elephant is like a great snake.''

The fourth touched a tusk.

``That cannot be right. An elephant is hard, smooth, and pointed. It is like a spear.''

The fifth held one of the ears as it moved in the air.

``You are describing something else entirely. An elephant is like a large fan.''

The sixth found the tail.

``All of you are making this far too complicated. An elephant is simply like a rope.''

Each person had touched the elephant.

Each person had reported honestly what they had found.

And their descriptions contradicted one another.

\section*{The quarrel}

At first, the disagreement was interesting.

Then it became personal.

The person at the flank could not understand how anyone could sincerely describe a wall as a snake.

The person at the trunk began to suspect that the others were incompetent.

The person at the tusk wondered whether some of them were lying.

Soon they were no longer talking about the elephant.

They were talking about one another.

``You never listen.''

``You cannot be trusted.''

``You always exaggerate.''

``You have no idea what you are doing.''

``There is no point asking you again.''

Their certainty grew at exactly the moment their access to one another's information began to disappear.

The village now had six descriptions of the elephant and less capacity to learn what an elephant actually was.

\section*{The easy solution}

The village council became tired of the argument.

One councillor proposed a simple solution.

``Choose one person,'' he said, ``and make that person's description official.''

The proposal was attractive.

It would end the dispute immediately.

If the person touching the flank were chosen, the village could write into its records:

\begin{quote}
\emph{An elephant is a wall.}
\end{quote}

The others could be instructed to stop confusing people.

There would be clarity.

There would be order.

There would also be a problem.

The official description might be wrong.

Worse, once disagreement itself became evidence of disloyalty to the official description, the village would have created a system that could no longer discover the mistake.

Another councillor suggested the opposite solution.

``Let everyone have their own truth. The person at the wall has a wall. The person at the trunk has a snake. We should stop trying to reconcile them.''

This protected everyone from being contradicted.

But it created a different problem.

There was still only one elephant.

And the elephant did not become six different animals simply because six people had encountered it differently.

If the village ever needed to understand the elephant well enough to care for it, avoid it, travel with it, or survive its movement through the village, private descriptions alone would not be enough.

They needed both freedom and contact.

\section*{The agreement}

So the six people tried something else.

They made an agreement.

The first rule was simple:

\begin{quote}
\textbf{No one would be forced to deny what they genuinely perceived.}
\end{quote}

If the leg felt like a tree, the person touching the leg was allowed to say so.

If the trunk felt like a snake, the person touching the trunk was allowed to say so.

No authority could reach into another person's experience and command:

\begin{quote}
``You did not feel what you felt.''
\end{quote}

But the agreement had a second rule:

\begin{quote}
\textbf{No person's perception would become unquestionable merely because it was sincerely held.}
\end{quote}

A person could say:

``This is what I feel.''

They could not conclude:

``Therefore nothing you feel can matter.''

Then they added more rules.

They would describe before accusing.

They would distinguish what they had directly felt from what they inferred.

They would ask questions when descriptions conflicted.

They would test claims where testing was possible.

They would change places.

They would return to parts of the elephant they thought they already understood.

They would allow new information to reopen old conclusions.

If someone had to be excluded from part of the inquiry because they were harming the process, that exclusion itself could be questioned later.

Most importantly, nobody could declare in advance:

\begin{quote}
``Nothing this person could ever discover will count.''
\end{quote}

The agreement did not require them to agree.

It required them to keep disagreement capable of carrying information.

\section*{Changing places}

The person who had touched the wall moved toward the trunk.

For the first time, he felt it twist.

``I understand why you called it a snake,'' he said.

The person who had held the trunk moved toward the leg.

``And I understand your tree.''

The person at the tusk explored the animal's head.

The person at the ear found the shoulder.

The person at the tail followed the animal's body forward.

Their earlier descriptions did not become false in the simple sense.

The flank really was broad.

The leg really was column-like.

The trunk really did move like a snake.

The tusk really was pointed.

Their mistake had been to confuse a true encounter with a complete model.

Gradually, a new description emerged.

Not because one person defeated the others.

Not because everyone surrendered their own experience.

And not because all descriptions were treated as equally accurate.

The shared model improved because differences could enter the relationship, be tested, and change what the group believed.

They began to understand something none of them had been able to establish alone:

the wall, tree, snake, spear, fan, and rope were connected.

They were parts of one animal.

\section*{When the elephant moved}

Then the elephant moved.

Suddenly the question was no longer philosophical.

Its weight shifted.

Its feet changed position.

Its trunk reached toward a basket.

People nearby had to decide where to stand.

Now the village's ability to combine partial information mattered.

The person near the front could feel the elephant turn before the person at the tail knew anything had changed.

The person at the leg could feel weight transfer.

The person at the flank could detect the movement of the whole body.

None had the entire picture.

But information could move between them.

Their interdependence had become a practical capacity.

The same relationship that helped them describe the elephant now helped them coordinate around it.

Their collective survival did not require anyone to become all-seeing.

It required a process in which partial observers could remain autonomous without becoming isolated, and connected without becoming subordinate to a single unquestionable observer.

\section*{What the agreement protected}

Years later, the village wrote down what it had learned.

It did not write:

\begin{quote}
\emph{Everyone is equally right.}
\end{quote}

That was clearly false.

It did not write:

\begin{quote}
\emph{Every claim must be discussed forever.}
\end{quote}

No village could function that way.

And it did not write:

\begin{quote}
\emph{No one may ever be excluded from anything.}
\end{quote}

Some boundaries were necessary.

Instead, the village protected two things.

First, it protected the person.

Each human being remained the first witness of what happened within their own experience: what they saw, felt, remembered, believed, feared, hoped, or understood.

No collective search for truth could justify taking ownership of another person's mind.

That was autonomy.

Second, it protected the person's place in the relationship.

Being fallible did not make a person disposable as a possible source of correction.

A person could be wrong.

A person could be ignored on a particular question for good reasons.

A claim could be rejected.

But the system itself had to remain capable of discovering that its rejection, exclusion, or current model was mistaken.

That was interdependence.

The villagers eventually expressed the symmetry this way:

\begin{quote}
\textbf{Protect the observer, because no collective has direct access to another person's experience.}

\textbf{Protect the relationship, because no observer has direct access to the whole truth.}
\end{quote}

\section*{The parable}

The elephant was reality.

The six people were human beings.

Their hands were perception, experience, memory, expertise, culture, measurement, and reasoning.

Their disagreement was not a defect in the existence of many minds.

It was information about the limits of each mind.

Their autonomy protected the integrity of what each person could genuinely observe and believe.

Their interdependence protected the possibility that those partial models could meet, challenge one another, and become something better.

Human rights are often understood as protections around the individual.

The parable suggests a complementary way to see them.

A human being is not only someone who must be protected from the power of others.

A human being is also a possible bearer of information that the rest of us do not yet have.

To protect the human is therefore to protect both dimensions:

\[
\boxed{\text{the autonomy to perceive, think, and speak from one's own position}}
\]

and

\[
\boxed{\text{the standing to remain part of a correctable shared search for reality}.}
\]

The first prevents the collective from swallowing the person.

The second prevents the collective from making the person epistemically nonexistent.

Neither is sufficient alone.

Autonomy without interdependence leaves six isolated truths that never become an elephant.

Interdependence without autonomy leaves one official description repeated by six mouths.

The agreement protects the space between those failures.

It allows a person to say:

\begin{quote}
``This is what I see.''
\end{quote}

and another to answer:

\begin{quote}
``I see something different.''
\end{quote}

without requiring either sentence to become an act of domination.

And so the deepest right in the story is not the right to be correct.

It is the right to remain a human participant in a process that can still become more correct.

\begin{center}
\bigskip
\textbf{None of them could see the whole elephant.}\\
\textbf{Together, under the agreement, they could keep learning it.}
\end{center}

\end{document}
